\section{Conclusion}
\label{sec:conclusion}
Tempo is the main nuisance of repetition counting. Rather than estimating it and
committing to one sampling rate, \method reads the video at several rates at once,
aligns the streams on one frame grid and lets a density decoder combine them. The
shared grid lets cached coarser streams serve as augmentation and as a second
test-time view, with exact input correspondence away from padding. With a frozen
encoder and a decoder of 3.9M parameters, \method reaches RepCount-A MAE 0.143,
below the reported values of the listed RGB-based counters under their respective
protocols. Rate-removal diagnostics show that coarse context matters more in
slow-tempo groups.

\paragraph{Limitations.} \method relies on a frozen encoder and on annotated
cycle boundaries for its density targets; Countix provides counts only, so its
targets place cycles uniformly. On UCFRep and Countix, counters that fine-tune
their encoder or pretrain on web-scale video--text data remain ahead
(\cref{app:other}), and adapting the encoder or adding a language model did not
close this gap consistently in our experiments (\cref{app:adapt,app:lm}).
Repetitions whose two half-strokes look alike remain the main source of
over-counting, and periodicity of frozen features alone does not disambiguate
them. The framework and main benchmark were selected after test inspection, and
external comparisons differ in training resources and metric conventions.
RepCount-A duplicate inputs and Countix shared source videos also leave benchmark
independence unestablished. These are historical original-split results;
independent confirmation after train-only correction remains open.
